% Options for packages loaded elsewhere
\PassOptionsToPackage{unicode}{hyperref}
\PassOptionsToPackage{hyphens}{url}
\documentclass[
  11pt,
  a4paper,
]{article}
\usepackage{xcolor}
\usepackage[margin=18mm]{geometry}
\usepackage{amsmath,amssymb}
\setcounter{secnumdepth}{-\maxdimen} % remove section numbering
\usepackage{iftex}
\ifPDFTeX
  \usepackage[T1]{fontenc}
  \usepackage[utf8]{inputenc}
  \usepackage{textcomp} % provide euro and other symbols
\else % if luatex or xetex
  \usepackage{unicode-math} % this also loads fontspec
  \defaultfontfeatures{Scale=MatchLowercase}
  \defaultfontfeatures[\rmfamily]{Ligatures=TeX,Scale=1}
\fi
\ifPDFTeX\else
  % xetex/luatex font selection
    \setmainfont[]{Noto Serif}
    \setsansfont[]{Noto Sans}
    \setmonofont[]{Noto Sans Mono}
  \setmathfont[]{Noto Sans Math}
\fi
% Use upquote if available, for straight quotes in verbatim environments
\IfFileExists{upquote.sty}{\usepackage{upquote}}{}
\IfFileExists{microtype.sty}{% use microtype if available
  \usepackage[]{microtype}
  \UseMicrotypeSet[protrusion]{basicmath} % disable protrusion for tt fonts
}{}
\makeatletter
\@ifundefined{KOMAClassName}{% if non-KOMA class
  \IfFileExists{parskip.sty}{%
    \usepackage{parskip}
  }{% else
    \setlength{\parindent}{0pt}
    \setlength{\parskip}{6pt plus 2pt minus 1pt}}
}{% if KOMA class
  \KOMAoptions{parskip=half}}
\makeatother
\ifLuaTeX
\usepackage[bidi=basic]{babel}
\else
\usepackage[bidi=default]{babel}
\fi
\babelprovide[main,import]{russian}
\ifPDFTeX
\else
\babelfont{rm}[]{Noto Serif}
\fi
% get rid of language-specific shorthands (see #6817):
\let\LanguageShortHands\languageshorthands
\def\languageshorthands#1{}
\setlength{\emergencystretch}{3em} % prevent overfull lines
\providecommand{\tightlist}{%
  \setlength{\itemsep}{0pt}\setlength{\parskip}{0pt}}
\usepackage{bookmark}
\IfFileExists{xurl.sty}{\usepackage{xurl}}{} % add URL line breaks if available
\urlstyle{same}
\hypersetup{
  pdftitle={Билеты по дискретным случайным процессам},
  pdflang={ru-RU},
  hidelinks,
  pdfcreator={LaTeX via pandoc}}

\title{Билеты по дискретным случайным процессам}
\author{}
\date{}

\begin{document}
\maketitle

\section{Билет 21. Закон больших чисел как частный случай центральной
предельной теоремы и как обоснование
аксиоматики}\label{ux431ux438ux43bux435ux442-21.-ux437ux430ux43aux43eux43d-ux431ux43eux43bux44cux448ux438ux445-ux447ux438ux441ux435ux43b-ux43aux430ux43a-ux447ux430ux441ux442ux43dux44bux439-ux441ux43bux443ux447ux430ux439-ux446ux435ux43dux442ux440ux430ux43bux44cux43dux43eux439-ux43fux440ux435ux434ux435ux43bux44cux43dux43eux439-ux442ux435ux43eux440ux435ux43cux44b-ux438-ux43aux430ux43a-ux43eux431ux43eux441ux43dux43eux432ux430ux43dux438ux435-ux430ux43aux441ux438ux43eux43cux430ux442ux438ux43aux438}

Источники: Ширяев 1, гл. I, §5-6; гл. III, §3; Ширяев 2, исторические
комментарии об аксиоматике Колмогорова.

\subsection{1. Короткий
ответ}\label{ux43aux43eux440ux43eux442ux43aux438ux439-ux43eux442ux432ux435ux442}

Закон больших чисел говорит, что среднее большого числа независимых
одинаково распределенных наблюдений сходится к математическому ожиданию:

\[
\overline X_n=\frac{X_1+\cdots+X_n}{n}\xrightarrow{P}m,
\qquad
m=\mathbb EX_1.
\]

Центральная предельная теорема уточняет масштаб случайных отклонений.
Если

\[
\mathbb EX_1=m,
\qquad
\operatorname{Var}X_1=\sigma^2\in(0,\infty),
\]

то

\[
\frac{X_1+\cdots+X_n-nm}{\sigma\sqrt n}
\Rightarrow
N(0,1).
\]

Из этой формулы видно, что

\[
X_1+\cdots+X_n-nm
\]

имеет типичный порядок \(\sqrt n\), поэтому после деления на \(n\)
отклонение среднего от \(m\) исчезает:

\[
\overline X_n-m
=
\frac{\sigma}{\sqrt n}
\frac{X_1+\cdots+X_n-nm}{\sigma\sqrt n}
\xrightarrow{P}0.
\]

Для индикаторов событий \(X_k=\mathbf 1_A^{(k)}\) это дает частотную
формулу

\[
\frac{N_n(A)}{n}\xrightarrow{P}P(A).
\]

Поэтому вероятность в аксиоматике Колмогорова не определяется как
частота, но закон больших чисел объясняет, почему вероятностная мера
согласуется с наблюдаемыми частотами в длинных сериях независимых
испытаний.

\subsection{2. Полный
ответ}\label{ux43fux43eux43bux43dux44bux439-ux43eux442ux432ux435ux442}

\subsubsection{2.1. Введение и основные
идеи}\label{ux432ux432ux435ux434ux435ux43dux438ux435-ux438-ux43eux441ux43dux43eux432ux43dux44bux435-ux438ux434ux435ux438}

В этом билете нужно связать три идеи:

\begin{itemize}
\tightlist
\item
  закон больших чисел;
\item
  центральную предельную теорему;
\item
  частотное обоснование вероятности как меры.
\end{itemize}

\subsubsection{2.2. Закон больших
чисел}\label{ux437ux430ux43aux43eux43d-ux431ux43eux43bux44cux448ux438ux445-ux447ux438ux441ux435ux43b}

Пусть

\[
X_1,X_2,\ldots
\]

\begin{itemize}
\tightlist
\item
  независимые одинаково распределенные случайные величины. Обозначим
\end{itemize}

\[
S_n=X_1+\cdots+X_n,
\qquad
\overline X_n=\frac{S_n}{n}.
\]

Если существует математическое ожидание

\[
\mathbb EX_1=m,
\]

то закон больших чисел утверждает, что выборочное среднее
\(\overline X_n\) при больших \(n\) близко к \(m\). В слабой форме это
записывается так:

\[
\overline X_n\xrightarrow{P}m,
\]

то есть для каждого \(\varepsilon>0\)

\[
P\{|\overline X_n-m|>\varepsilon\}\to0.
\]

Интуитивно: отдельное наблюдение может сильно колебаться, но при
усреднении независимые случайные отклонения взаимно компенсируются.

\subsubsection{2.3. Центральная предельная
теорема}\label{ux446ux435ux43dux442ux440ux430ux43bux44cux43dux430ux44f-ux43fux440ux435ux434ux435ux43bux44cux43dux430ux44f-ux442ux435ux43eux440ux435ux43cux430}

Центральная предельная теорема дает более точное утверждение. Если,
кроме среднего, существует конечная ненулевая дисперсия

\[
\operatorname{Var}X_1=\sigma^2\in(0,\infty),
\]

то

\[
\frac{S_n-nm}{\sigma\sqrt n}
\Rightarrow
N(0,1).
\]

Здесь \(\Rightarrow\) означает сходимость по распределению.
Эквивалентно, для всех точек непрерывности функции распределения
стандартного нормального закона:

\[
P\left\{
\frac{S_n-nm}{\sigma\sqrt n}\le x
\right\}
\to
\Phi(x),
\]

где

\[
\Phi(x)=\frac{1}{\sqrt{2\pi}}\int_{-\infty}^x e^{-u^2/2}\,du.
\]

\subsubsection{2.4. Вывод ЗБЧ из
ЦПТ}\label{ux432ux44bux432ux43eux434-ux437ux431ux447-ux438ux437-ux446ux43fux442}

Теперь покажем, почему ЗБЧ можно получить из ЦПТ. Обозначим

\[
Z_n=
\frac{S_n-nm}{\sigma\sqrt n}.
\]

По ЦПТ

\[
Z_n\Rightarrow Z,
\qquad
Z\sim N(0,1).
\]

Последовательность, сходящаяся по распределению к настоящей случайной
величине, ограничена по вероятности: ее значения не уходят на
бесконечность с положительной вероятностью. Поэтому

\[
Z_n=O_P(1).
\]

Тогда

\[
\overline X_n-m
=
\frac{S_n-nm}{n}
=
\frac{\sigma}{\sqrt n}Z_n
\xrightarrow{P}0.
\]

Более подробно: для любого \(\varepsilon>0\)

\[
P\{|\overline X_n-m|>\varepsilon\}
=
P\left\{|Z_n|>\frac{\varepsilon\sqrt n}{\sigma}\right\}.
\]

Правая граница уходит в бесконечность, а \(Z_n\) остается ограниченной
по вероятности, значит эта вероятность стремится к нулю. Это и есть
слабый закон больших чисел.

\subsubsection{2.5. Условия применимости
теорем}\label{ux443ux441ux43bux43eux432ux438ux44f-ux43fux440ux438ux43cux435ux43dux438ux43cux43eux441ux442ux438-ux442ux435ux43eux440ux435ux43c}

Важно сказать на экзамене: не всякий закон больших чисел является
буквально частным случаем ЦПТ. ЦПТ требует конечной ненулевой дисперсии,
а слабый закон больших чисел может быть верен при более слабых условиях,
например для iid-величин с конечным первым моментом. Поэтому корректная
фраза такая: \textbf{при условиях ЦПТ закон больших чисел следует из нее
как более грубое утверждение}.

\subsubsection{2.6. Схема Бернулли и частотное
обоснование}\label{ux441ux445ux435ux43cux430-ux431ux435ux440ux43dux443ux43bux43bux438-ux438-ux447ux430ux441ux442ux43eux442ux43dux43eux435-ux43eux431ux43eux441ux43dux43eux432ux430ux43dux438ux435}

Исторически и методически главный пример - схема Бернулли. Пусть
проводится последовательность независимых испытаний, и событие \(A\)
происходит в каждом испытании с вероятностью

\[
p=P(A).
\]

Положим

\[
X_k=\mathbf 1_{\{A\text{ произошло в }k\text{-м испытании}\}}.
\]

Тогда

\[
P\{X_k=1\}=p,
\qquad
P\{X_k=0\}=q=1-p,
\]

и

\[
\mathbb EX_k=p,
\qquad
\operatorname{Var}X_k=pq.
\]

Число появлений события \(A\) за \(n\) испытаний равно

\[
N_n(A)=X_1+\cdots+X_n.
\]

Относительная частота равна

\[
\nu_n(A)=\frac{N_n(A)}{n}.
\]

Закон больших чисел Бернулли утверждает:

\[
\frac{N_n(A)}{n}\xrightarrow{P}p.
\]

То есть для каждого \(\varepsilon>0\)

\[
P\left\{
\left|\frac{N_n(A)}{n}-p\right|>\varepsilon
\right\}
\to0.
\]

Если \(0<p<1\), то из ЦПТ для индикаторов получаем более точную формулу:

\[
\frac{N_n(A)-np}{\sqrt{npq}}
\Rightarrow
N(0,1).
\]

Отсюда снова следует

\[
\frac{N_n(A)}{n}-p
=
\sqrt{\frac{pq}{n}}\,
\frac{N_n(A)-np}{\sqrt{npq}}
\xrightarrow{P}0.
\]

\subsubsection{2.7. Связь с аксиоматикой
Колмогорова}\label{ux441ux432ux44fux437ux44c-ux441-ux430ux43aux441ux438ux43eux43cux430ux442ux438ux43aux43eux439-ux43aux43eux43bux43cux43eux433ux43eux440ux43eux432ux430}

Именно этот пример объясняет связь с аксиоматикой. В аксиоматике
Колмогорова вероятность - это не предел частоты по определению, а мера

\[
P:\mathcal F\to[0,1]
\]

на \(\sigma\)-алгебре событий, такая что

\[
P(\Omega)=1
\]

и для попарно несовместных событий \(A_1,A_2,\ldots\)

\[
P\left(\bigcup_{k=1}^{\infty}A_k\right)
=
\sum_{k=1}^{\infty}P(A_k).
\]

Такой подход математически удобен: он позволяет строить случайные
величины, распределения, интегралы, условные ожидания и случайные
процессы. Но возникает естественный вопрос: почему это число \(P(A)\)
имеет отношение к эксперименту?

Ответ дает закон больших чисел: если событие \(A\) наблюдается в длинной
серии независимых повторений одного и того же опыта, то относительная
частота события будет близка к \(P(A)\) с вероятностью, стремящейся к
единице. Поэтому мера \(P\) получает эмпирическую интерпретацию через
устойчивость частот.

\subsubsection{2.8. Эмпирическая интерпретация и применимость
модели}\label{ux44dux43cux43fux438ux440ux438ux447ux435ux441ux43aux430ux44f-ux438ux43dux442ux435ux440ux43fux440ux435ux442ux430ux446ux438ux44f-ux438-ux43fux440ux438ux43cux435ux43dux438ux43cux43eux441ux442ux44c-ux43cux43eux434ux435ux43bux438}

Тонкость: закон больших чисел не доказывает аксиомы вероятности из
физики. Он показывает, что если мы приняли колмогоровскую модель и
независимость повторных испытаний, то модель предсказывает частотную
устойчивость. Поэтому ЗБЧ является не логическим выводом аксиом из
частот, а математическим обоснованием применимости аксиоматики к
повторяемым экспериментам.

\subsection{3. Основные
определения}\label{ux43eux441ux43dux43eux432ux43dux44bux435-ux43eux43fux440ux435ux434ux435ux43bux435ux43dux438ux44f}

\subsubsection{Сходимость по
вероятности}\label{ux441ux445ux43eux434ux438ux43cux43eux441ux442ux44c-ux43fux43e-ux432ux435ux440ux43eux44fux442ux43dux43eux441ux442ux438}

Последовательность случайных величин \(Y_n\) сходится по вероятности к
\(Y\), если для каждого \(\varepsilon>0\)

\[
P\{|Y_n-Y|>\varepsilon\}\to0.
\]

Обозначение:

\[
Y_n\xrightarrow{P}Y.
\]

\subsubsection{Сходимость по
распределению}\label{ux441ux445ux43eux434ux438ux43cux43eux441ux442ux44c-ux43fux43e-ux440ux430ux441ux43fux440ux435ux434ux435ux43bux435ux43dux438ux44e}

Последовательность \(Y_n\) сходится по распределению к \(Y\), если

\[
F_{Y_n}(x)\to F_Y(x)
\]

во всех точках непрерывности \(F_Y\). Обозначение:

\[
Y_n\Rightarrow Y.
\]

Эквивалентно, для всех ограниченных непрерывных функций \(f\):

\[
\mathbb Ef(Y_n)\to\mathbb Ef(Y).
\]

Эта формулировка связана с Билетом 20, где распределения рассматриваются
как функционалы на тестовых функциях.

\subsubsection{Ограниченность по
вероятности}\label{ux43eux433ux440ux430ux43dux438ux447ux435ux43dux43dux43eux441ux442ux44c-ux43fux43e-ux432ux435ux440ux43eux44fux442ux43dux43eux441ux442ux438}

Последовательность \(Y_n\) называется ограниченной по вероятности, если
для каждого \(\delta>0\) существует \(M>0\) такое, что при всех
достаточно больших \(n\)

\[
P\{|Y_n|>M\}<\delta.
\]

Пишут:

\[
Y_n=O_P(1).
\]

Если \(Y_n\Rightarrow Y\), то \(Y_n=O_P(1)\).

\subsubsection{Слабый закон больших
чисел}\label{ux441ux43bux430ux431ux44bux439-ux437ux430ux43aux43eux43d-ux431ux43eux43bux44cux448ux438ux445-ux447ux438ux441ux435ux43b}

Для независимых одинаково распределенных величин \(X_1,X_2,\ldots\) с
\(\mathbb EX_1=m\) слабый ЗБЧ утверждает

\[
\frac{X_1+\cdots+X_n}{n}\xrightarrow{P}m
\]

при соответствующих условиях. Для iid-величин достаточно
\(\mathbb E|X_1|<\infty\).

В частном случае конечной дисперсии это легко доказывается неравенством
Чебышева.

\subsubsection{Усиленный закон больших
чисел}\label{ux443ux441ux438ux43bux435ux43dux43dux44bux439-ux437ux430ux43aux43eux43d-ux431ux43eux43bux44cux448ux438ux445-ux447ux438ux441ux435ux43b}

Усиленный ЗБЧ утверждает почти наверное сходимость:

\[
\frac{X_1+\cdots+X_n}{n}\to m
\quad\text{п.н.}
\]

Это более сильное утверждение, чем сходимость по вероятности.

\subsubsection{Центральная предельная
теорема}\label{ux446ux435ux43dux442ux440ux430ux43bux44cux43dux430ux44f-ux43fux440ux435ux434ux435ux43bux44cux43dux430ux44f-ux442ux435ux43eux440ux435ux43cux430-1}

Если \(X_1,X_2,\ldots\) - iid-величины,

\[
\mathbb EX_1=m,
\qquad
\operatorname{Var}X_1=\sigma^2\in(0,\infty),
\]

то

\[
\frac{S_n-nm}{\sigma\sqrt n}
\Rightarrow
N(0,1).
\]

\subsubsection{Схема
Бернулли}\label{ux441ux445ux435ux43cux430-ux431ux435ux440ux43dux443ux43bux43bux438}

Это последовательность независимых испытаний с двумя исходами: успех с
вероятностью \(p\) и неудача с вероятностью \(q=1-p\). Если \(X_k\) -
индикатор успеха, то

\[
S_n=X_1+\cdots+X_n
\]

имеет биномиальное распределение:

\[
P\{S_n=k\}=C_n^k p^kq^{n-k}.
\]

\subsubsection{Относительная
частота}\label{ux43eux442ux43dux43eux441ux438ux442ux435ux43bux44cux43dux430ux44f-ux447ux430ux441ux442ux43eux442ux430}

Если событие \(A\) произошло \(N_n(A)\) раз в \(n\) испытаниях, то

\[
\nu_n(A)=\frac{N_n(A)}{n}
\]

называется относительной частотой события \(A\).

\subsection{4. Основные теоремы и
свойства}\label{ux43eux441ux43dux43eux432ux43dux44bux435-ux442ux435ux43eux440ux435ux43cux44b-ux438-ux441ux432ux43eux439ux441ux442ux432ux430}

\subsubsection{1. Неравенство
Чебышева}\label{ux43dux435ux440ux430ux432ux435ux43dux441ux442ux432ux43e-ux447ux435ux431ux44bux448ux435ux432ux430}

Если случайная величина \(Y\) имеет конечную дисперсию, то для любого
\(\varepsilon>0\)

\[
P\{|Y-\mathbb EY|\ge\varepsilon\}
\le
\frac{\operatorname{Var}Y}{\varepsilon^2}.
\]

Это базовый инструмент доказательства слабого ЗБЧ при конечной
дисперсии.

\subsubsection{2. Слабый ЗБЧ при конечной
дисперсии}\label{ux441ux43bux430ux431ux44bux439-ux437ux431ux447-ux43fux440ux438-ux43aux43eux43dux435ux447ux43dux43eux439-ux434ux438ux441ux43fux435ux440ux441ux438ux438}

Пусть \(X_1,X_2,\ldots\) независимы, одинаково распределены,

\[
\mathbb EX_1=m,
\qquad
\operatorname{Var}X_1=\sigma^2<\infty.
\]

Тогда

\[
\overline X_n\xrightarrow{P}m.
\]

Доказательство:

\[
\mathbb E\overline X_n=m,
\]

а по независимости

\[
\operatorname{Var}\overline X_n
=
\operatorname{Var}\left(\frac{S_n}{n}\right)
=
\frac{1}{n^2}\operatorname{Var}S_n
=
\frac{1}{n^2}n\sigma^2
=
\frac{\sigma^2}{n}.
\]

По неравенству Чебышева:

\[
P\{|\overline X_n-m|>\varepsilon\}
\le
\frac{\sigma^2}{n\varepsilon^2}
\to0.
\]

\subsubsection{3. ЦПТ для
iid-величин}\label{ux446ux43fux442-ux434ux43bux44f-iid-ux432ux435ux43bux438ux447ux438ux43d}

При условиях

\[
\mathbb EX_1=m,
\qquad
\operatorname{Var}X_1=\sigma^2\in(0,\infty)
\]

верна сходимость

\[
\frac{S_n-nm}{\sigma\sqrt n}
\Rightarrow
N(0,1).
\]

Смысл: сумма \(S_n\) сосредоточена около \(nm\), а ее случайные
флуктуации имеют масштаб \(\sqrt n\).

\subsubsection{4. ЗБЧ как следствие
ЦПТ}\label{ux437ux431ux447-ux43aux430ux43a-ux441ux43bux435ux434ux441ux442ux432ux438ux435-ux446ux43fux442}

Если выполнены условия ЦПТ, то

\[
\overline X_n-m
=
\frac{\sigma}{\sqrt n}
\frac{S_n-nm}{\sigma\sqrt n}.
\]

Второй множитель сходится по распределению к \(N(0,1)\), значит он
ограничен по вероятности. Первый множитель стремится к нулю.
Следовательно,

\[
\overline X_n-m\xrightarrow{P}0.
\]

Это и есть слабый ЗБЧ.

\subsubsection{5. Закон больших чисел
Бернулли}\label{ux437ux430ux43aux43eux43d-ux431ux43eux43bux44cux448ux438ux445-ux447ux438ux441ux435ux43b-ux431ux435ux440ux43dux443ux43bux43bux438}

В схеме Бернулли

\[
S_n\sim\operatorname{Bin}(n,p).
\]

Тогда

\[
\frac{S_n}{n}\xrightarrow{P}p.
\]

Через неравенство Чебышева:

\[
\operatorname{Var}\frac{S_n}{n}
=
\frac{pq}{n},
\]

поэтому

\[
P\left\{
\left|\frac{S_n}{n}-p\right|>\varepsilon
\right\}
\le
\frac{pq}{n\varepsilon^2}
\to0.
\]

Через ЦПТ:

\[
\frac{S_n-np}{\sqrt{npq}}
\Rightarrow
N(0,1),
\]

откуда

\[
\frac{S_n}{n}-p
=
\sqrt{\frac{pq}{n}}
\frac{S_n-np}{\sqrt{npq}}
\xrightarrow{P}0.
\]

\subsubsection{6. Частотная интерпретация
вероятности}\label{ux447ux430ux441ux442ux43eux442ux43dux430ux44f-ux438ux43dux442ux435ux440ux43fux440ux435ux442ux430ux446ux438ux44f-ux432ux435ux440ux43eux44fux442ux43dux43eux441ux442ux438}

Если \(A\) - событие с вероятностью \(P(A)=p\), а \(\nu_n(A)\) - частота
его появления в \(n\) независимых повторениях опыта, то

\[
\nu_n(A)\xrightarrow{P}P(A).
\]

Это объясняет, почему вероятность можно практически оценивать частотой:

\[
P(A)\approx \frac{N_n(A)}{n}
\]

при большом \(n\).

\subsection{5. Примеры}\label{ux43fux440ux438ux43cux435ux440ux44b}

\subsubsection{Пример 1. Подбрасывание
монеты}\label{ux43fux440ux438ux43cux435ux440-1.-ux43fux43eux434ux431ux440ux430ux441ux44bux432ux430ux43dux438ux435-ux43cux43eux43dux435ux442ux44b}

Пусть \(A\) - выпадение орла, \(P(A)=p\). В \(n\) независимых
подбрасываниях

\[
N_n(A)\sim\operatorname{Bin}(n,p).
\]

Тогда

\[
\frac{N_n(A)}{n}\xrightarrow{P}p.
\]

Если монета правильная, \(p=1/2\), и

\[
\frac{N_n(A)-n/2}{\sqrt{n}/2}
\Rightarrow
N(0,1).
\]

ЗБЧ говорит, что доля орлов близка к \(1/2\), а ЦПТ говорит, какого
порядка типичное отклонение:

\[
N_n(A)-\frac n2
\]

обычно имеет порядок \(\sqrt n\), а не порядок \(n\).

\subsubsection{Пример 2. Среднее
измерений}\label{ux43fux440ux438ux43cux435ux440-2.-ux441ux440ux435ux434ux43dux435ux435-ux438ux437ux43cux435ux440ux435ux43dux438ux439}

Пусть измеряется величина \(m\), а результат \(k\)-го измерения:

\[
X_k=m+\xi_k,
\]

где ошибки \(\xi_k\) независимы,

\[
\mathbb E\xi_k=0,
\qquad
\operatorname{Var}\xi_k=\sigma^2.
\]

Тогда среднее измерений

\[
\overline X_n=m+\frac{\xi_1+\cdots+\xi_n}{n}
\]

сходится к \(m\) по вероятности. ЦПТ уточняет:

\[
\frac{\sqrt n(\overline X_n-m)}{\sigma}
\Rightarrow
N(0,1).
\]

Поэтому ошибка среднего имеет порядок

\[
\frac{\sigma}{\sqrt n}.
\]

\subsubsection{Пример 3. Оценка вероятности
события}\label{ux43fux440ux438ux43cux435ux440-3.-ux43eux446ux435ux43dux43aux430-ux432ux435ux440ux43eux44fux442ux43dux43eux441ux442ux438-ux441ux43eux431ux44bux442ux438ux44f}

Пусть нужно оценить вероятность отказа устройства \(A\). Проводим \(n\)
независимых испытаний и ставим

\[
X_k=\mathbf 1_{\{A\text{ произошло в }k\text{-м испытании}\}}.
\]

Тогда

\[
\widehat p_n=\frac{1}{n}\sum_{k=1}^nX_k
\]

\begin{itemize}
\tightlist
\item
  естественная оценка \(p=P(A)\). По ЗБЧ
\end{itemize}

\[
\widehat p_n\xrightarrow{P}p.
\]

Если \(0<p<1\), то по ЦПТ

\[
\frac{\widehat p_n-p}{\sqrt{p(1-p)/n}}
\Rightarrow
N(0,1).
\]

Это основа нормальных доверительных интервалов для вероятности.

\subsubsection{Пример 4. Случайное
блуждание}\label{ux43fux440ux438ux43cux435ux440-4.-ux441ux43bux443ux447ux430ux439ux43dux43eux435-ux431ux43bux443ux436ux434ux430ux43dux438ux435}

Пусть

\[
S_n=\xi_1+\cdots+\xi_n,
\]

где \(\xi_k\) - iid-шаги,

\[
\mathbb E\xi_k=a,
\qquad
\operatorname{Var}\xi_k=\sigma^2.
\]

Тогда

\[
\frac{S_n}{n}\xrightarrow{P}a,
\]

а при \(\sigma^2>0\)

\[
\frac{S_n-na}{\sigma\sqrt n}
\Rightarrow
N(0,1).
\]

Это связывает билет с Билетом 19: траектория блуждания в среднем идет с
линейной скоростью \(a\), а случайные отклонения от прямой имеют масштаб
\(\sqrt n\).

\subsection{6. Схемы доказательств /
обоснований}\label{ux441ux445ux435ux43cux44b-ux434ux43eux43aux430ux437ux430ux442ux435ux43bux44cux441ux442ux432-ux43eux431ux43eux441ux43dux43eux432ux430ux43dux438ux439}

\textbf{Доказательство слабого ЗБЧ через Чебышева.}

\begin{enumerate}
\def\labelenumi{\arabic{enumi}.}
\tightlist
\item
  Берем
\end{enumerate}

\[
\overline X_n=\frac{S_n}{n}.
\]

\begin{enumerate}
\def\labelenumi{\arabic{enumi}.}
\setcounter{enumi}{1}
\tightlist
\item
  Считаем матожидание:
\end{enumerate}

\[
\mathbb E\overline X_n=m.
\]

\begin{enumerate}
\def\labelenumi{\arabic{enumi}.}
\setcounter{enumi}{2}
\tightlist
\item
  Считаем дисперсию, используя независимость:
\end{enumerate}

\[
\operatorname{Var}\overline X_n
=
\frac{\sigma^2}{n}.
\]

\begin{enumerate}
\def\labelenumi{\arabic{enumi}.}
\setcounter{enumi}{3}
\tightlist
\item
  Применяем Чебышева:
\end{enumerate}

\[
P\{|\overline X_n-m|>\varepsilon\}
\le
\frac{\sigma^2}{n\varepsilon^2}.
\]

\begin{enumerate}
\def\labelenumi{\arabic{enumi}.}
\setcounter{enumi}{4}
\tightlist
\item
  Правая часть стремится к нулю, значит
\end{enumerate}

\[
\overline X_n\xrightarrow{P}m.
\]

\textbf{Доказательство ЗБЧ как следствия ЦПТ.}

\begin{enumerate}
\def\labelenumi{\arabic{enumi}.}
\tightlist
\item
  По ЦПТ:
\end{enumerate}

\[
Z_n=\frac{S_n-nm}{\sigma\sqrt n}
\Rightarrow
Z\sim N(0,1).
\]

\begin{enumerate}
\def\labelenumi{\arabic{enumi}.}
\setcounter{enumi}{1}
\item
  Значит \(Z_n=O_P(1)\).
\item
  Переписываем отклонение среднего:
\end{enumerate}

\[
\overline X_n-m
=
\frac{\sigma}{\sqrt n}Z_n.
\]

\begin{enumerate}
\def\labelenumi{\arabic{enumi}.}
\setcounter{enumi}{3}
\tightlist
\item
  Произведение ограниченной по вероятности последовательности на
  числовую последовательность, стремящуюся к нулю, сходится к нулю по
  вероятности:
\end{enumerate}

\[
\frac{\sigma}{\sqrt n}Z_n\xrightarrow{P}0.
\]

\begin{enumerate}
\def\labelenumi{\arabic{enumi}.}
\setcounter{enumi}{4}
\tightlist
\item
  Следовательно,
\end{enumerate}

\[
\overline X_n\xrightarrow{P}m.
\]

\textbf{Доказательство частотной формулы.}

\begin{enumerate}
\def\labelenumi{\arabic{enumi}.}
\tightlist
\item
  Для события \(A\) вводим индикаторы:
\end{enumerate}

\[
X_k=\mathbf 1_{\{A\text{ произошло в }k\text{-м испытании}\}}.
\]

\begin{enumerate}
\def\labelenumi{\arabic{enumi}.}
\setcounter{enumi}{1}
\tightlist
\item
  Тогда
\end{enumerate}

\[
\mathbb EX_k=P(A)=p.
\]

\begin{enumerate}
\def\labelenumi{\arabic{enumi}.}
\setcounter{enumi}{2}
\tightlist
\item
  Частота равна среднему индикаторов:
\end{enumerate}

\[
\nu_n(A)=\frac{X_1+\cdots+X_n}{n}.
\]

\begin{enumerate}
\def\labelenumi{\arabic{enumi}.}
\setcounter{enumi}{3}
\tightlist
\item
  По ЗБЧ:
\end{enumerate}

\[
\nu_n(A)\xrightarrow{P}p=P(A).
\]

Это и есть математическая форма утверждения: ``вероятность проявляется
как устойчивая частота в длинной серии независимых испытаний''.

\subsection{7. Что написать на
доске}\label{ux447ux442ux43e-ux43dux430ux43fux438ux441ux430ux442ux44c-ux43dux430-ux434ux43eux441ux43aux435}

Минимальный набор формул:

\[
S_n=X_1+\cdots+X_n,
\qquad
\overline X_n=\frac{S_n}{n}.
\]

ЗБЧ:

\[
\overline X_n\xrightarrow{P}m,
\qquad
m=\mathbb EX_1.
\]

ЦПТ:

\[
\frac{S_n-nm}{\sigma\sqrt n}
\Rightarrow
N(0,1).
\]

Вывод ЗБЧ из ЦПТ:

\[
\overline X_n-m
=
\frac{\sigma}{\sqrt n}
\frac{S_n-nm}{\sigma\sqrt n}
\xrightarrow{P}0.
\]

Частоты:

\[
X_k=\mathbf 1_A^{(k)},
\qquad
\nu_n(A)=\frac{1}{n}\sum_{k=1}^nX_k,
\qquad
\nu_n(A)\xrightarrow{P}P(A).
\]

Фраза про аксиоматику:

\[
P(A)\ \text{не определяется как частота, но ЗБЧ объясняет частотный смысл }P(A).
\]

\subsection{8. Типичные дополнительные
вопросы}\label{ux442ux438ux43fux438ux447ux43dux44bux435-ux434ux43eux43fux43eux43bux43dux438ux442ux435ux43bux44cux43dux44bux435-ux432ux43eux43fux440ux43eux441ux44b}

\textbf{В каком смысле ЗБЧ является частным случаем ЦПТ?}

При условиях ЦПТ нормированное отклонение

\[
\frac{S_n-nm}{\sigma\sqrt n}
\]

имеет невырожденный предельный закон, значит отклонение \(S_n-nm\) имеет
порядок \(\sqrt n\). После деления на \(n\) получаем отклонение среднего
порядка \(1/\sqrt n\), которое стремится к нулю по вероятности.

\textbf{Почему нельзя сказать, что любой ЗБЧ следует из ЦПТ?}

Потому что ЦПТ требует конечной ненулевой дисперсии, а некоторые формы
ЗБЧ верны при более слабых условиях. Например, для iid-величин слабый
ЗБЧ верен при \(\mathbb E|X_1|<\infty\), даже если дисперсия бесконечна.

\textbf{Что сильнее: сходимость по вероятности или по распределению?}

Сходимость по вероятности к константе влечет сходимость по распределению
к этой константе. Обратно в общем случае неверно, но если

\[
Y_n\Rightarrow c
\]

к константе \(c\), то

\[
Y_n\xrightarrow{P}c.
\]

В выводе ЗБЧ из ЦПТ мы используем не это напрямую, а факт, что
нормированная сумма ограничена по вероятности.

\textbf{Зачем в ЦПТ нормировка именно на \(\sqrt n\)?}

Для независимых одинаково распределенных величин дисперсии суммируются:

\[
\operatorname{Var}S_n=n\sigma^2.
\]

Стандартное отклонение суммы равно

\[
\sigma\sqrt n.
\]

Поэтому именно деление на \(\sigma\sqrt n\) дает невырожденный
предельный закон.

\textbf{Как связаны ЗБЧ и оценивание вероятности?}

Если \(X_k=\mathbf 1_A^{(k)}\), то

\[
\mathbb EX_k=P(A).
\]

Среднее индикаторов равно частоте события. По ЗБЧ частота сходится к
\(P(A)\), поэтому частота является состоятельной оценкой вероятности.

\textbf{Что именно обосновывает ЗБЧ в аксиоматике Колмогорова?}

Он обосновывает интерпретацию вероятности как идеального значения, около
которого стабилизируются частоты в независимых повторениях. Но сами
аксиомы - \(\sigma\)-алгебра, мера, счетная аддитивность - остаются
математическими предположениями модели.

\textbf{Почему независимость важна?}

Без независимости среднее может не стабилизироваться около
математического ожидания. Например, если все \(X_k\) равны одной и той
же случайной величине \(X\), то

\[
\overline X_n=X
\]

для всех \(n\), и среднее не обязано сходиться к \(\mathbb EX\) по
вероятности.

\textbf{Что будет при \(\sigma^2=0\)?}

Тогда \(X_1=m\) почти наверное, и

\[
\overline X_n=m
\]

почти наверное для всех \(n\). ЗБЧ тривиален, а стандартная формула ЦПТ
с делением на \(\sigma\) неприменима.

\subsection{9. Типичные
ошибки}\label{ux442ux438ux43fux438ux447ux43dux44bux435-ux43eux448ux438ux431ux43aux438}

\begin{itemize}
\tightlist
\item
  Говорить, что частота ``равна'' вероятности при большом \(n\).
  Правильно: частота близка к вероятности с большой вероятностью.
\item
  Забывать вид сходимости: в слабом ЗБЧ это сходимость по вероятности,
  не поточечная сходимость для каждого исхода.
\item
  Называть ЗБЧ частным случаем ЦПТ без оговорки о конечной дисперсии.
\item
  Путать масштаб ЦПТ и ЗБЧ: сумма отклоняется на \(\sqrt n\), среднее
  отклоняется на \(1/\sqrt n\).
\item
  Забывать независимость или хотя бы условия слабой зависимости.
\item
  Для схемы Бернулли писать дисперсию частоты как \(npq\) вместо
  \(pq/n\).
\item
  Считать, что ЗБЧ доказывает существование вероятности как предела
  частот. В колмогоровской теории вероятность задается мерой, а ЗБЧ
  объясняет частотную устойчивость внутри модели.
\end{itemize}

\subsection{10. Связи с другими
билетами}\label{ux441ux432ux44fux437ux438-ux441-ux434ux440ux443ux433ux438ux43cux438-ux431ux438ux43bux435ux442ux430ux43cux438}

\begin{itemize}
\tightlist
\item
  Билет 02: вероятностное пространство Колмогорова и счетная
  аддитивность меры.
\item
  Билет 08: характеристические функции используются в доказательстве
  ЦПТ.
\item
  Билет 09: независимость нужна для произведения характеристических
  функций и сложения дисперсий.
\item
  Билет 11: математическое ожидание и дисперсия выражаются как
  интегралы.
\item
  Билет 13: моменты связаны с производными характеристической функции.
\item
  Билет 16: последовательность наблюдений можно рассматривать как
  дискретный случайный процесс.
\item
  Билет 17: эргодичность обобщает идею сходимости временных средних.
\item
  Билет 19: для случайных блужданий ЗБЧ и ЦПТ описывают асимптотику
  суммы шагов.
\item
  Билет 20: слабая сходимость распределений формулируется через тестовые
  функции.
\item
  Билет 24: достаточные условия эргодичности дают аналоги ЗБЧ для
  зависимых стационарных последовательностей.
\end{itemize}

\subsection{11. Минимум на
удовлетворительно}\label{ux43cux438ux43dux438ux43cux443ux43c-ux43dux430-ux443ux434ux43eux432ux43bux435ux442ux432ux43eux440ux438ux442ux435ux43bux44cux43dux43e}

Нужно сказать:

\begin{enumerate}
\def\labelenumi{\arabic{enumi}.}
\tightlist
\item
  ЗБЧ:
\end{enumerate}

\[
\frac{X_1+\cdots+X_n}{n}\xrightarrow{P}\mathbb EX_1.
\]

\begin{enumerate}
\def\labelenumi{\arabic{enumi}.}
\setcounter{enumi}{1}
\tightlist
\item
  ЦПТ:
\end{enumerate}

\[
\frac{X_1+\cdots+X_n-nm}{\sigma\sqrt n}
\Rightarrow
N(0,1).
\]

\begin{enumerate}
\def\labelenumi{\arabic{enumi}.}
\setcounter{enumi}{2}
\item
  Из ЦПТ следует ЗБЧ, потому что отклонение суммы от \(nm\) имеет
  порядок \(\sqrt n\), а после деления на \(n\) становится малым.
\item
  Для индикаторов событий:
\end{enumerate}

\[
\frac{N_n(A)}{n}\xrightarrow{P}P(A).
\]

\begin{enumerate}
\def\labelenumi{\arabic{enumi}.}
\setcounter{enumi}{4}
\tightlist
\item
  Это объясняет частотный смысл вероятности.
\end{enumerate}

\subsection{12. Минимум на
хорошо}\label{ux43cux438ux43dux438ux43cux443ux43c-ux43dux430-ux445ux43eux440ux43eux448ux43e}

Нужно дополнительно:

\begin{itemize}
\tightlist
\item
  строго определить сходимость по вероятности и сходимость по
  распределению;
\item
  назвать условия ЦПТ: iid, конечное среднее, конечная ненулевая
  дисперсия;
\item
  аккуратно вывести
\end{itemize}

\[
\overline X_n-m
=
\frac{\sigma}{\sqrt n}Z_n;
\]

\begin{itemize}
\tightlist
\item
  разобрать схему Бернулли;
\item
  сказать, что ЗБЧ шире ЦПТ и не всегда требует конечной дисперсии.
\end{itemize}

\subsection{13. Что добавить для
отлично}\label{ux447ux442ux43e-ux434ux43eux431ux430ux432ux438ux442ux44c-ux434ux43bux44f-ux43eux442ux43bux438ux447ux43dux43e}

Для сильного ответа стоит добавить:

\begin{itemize}
\tightlist
\item
  доказательство слабого ЗБЧ через Чебышева;
\item
  объяснение ограниченности по вероятности при сходимости по
  распределению;
\item
  различие между слабым и усиленным ЗБЧ;
\item
  фразу о том, что аксиоматика не выводится из частот, а получает через
  ЗБЧ экспериментальную интерпретацию;
\item
  связь с эргодичностью: в более общих зависимых процессах роль ЗБЧ
  играют эргодические теоремы.
\end{itemize}

\subsection{14. Устная версия ответа на 5
минут}\label{ux443ux441ux442ux43dux430ux44f-ux432ux435ux440ux441ux438ux44f-ux43eux442ux432ux435ux442ux430-ux43dux430-5-ux43cux438ux43dux443ux442}

Закон больших чисел и центральная предельная теорема описывают
асимптотику сумм независимых одинаково распределенных случайных величин.
Пусть

\[
S_n=X_1+\cdots+X_n,
\qquad
\mathbb EX_1=m.
\]

Слабый закон больших чисел утверждает:

\[
\frac{S_n}{n}\xrightarrow{P}m.
\]

То есть для каждого \(\varepsilon>0\)

\[
P\left\{\left|\frac{S_n}{n}-m\right|>\varepsilon\right\}\to0.
\]

Центральная предельная теорема требует дополнительно конечной ненулевой
дисперсии

\[
\operatorname{Var}X_1=\sigma^2\in(0,\infty)
\]

и утверждает:

\[
\frac{S_n-nm}{\sigma\sqrt n}
\Rightarrow
N(0,1).
\]

Из этой формулы сразу видно, что отклонение суммы от \(nm\) имеет
масштаб \(\sqrt n\). Поэтому

\[
\frac{S_n}{n}-m
=
\frac{\sigma}{\sqrt n}
\frac{S_n-nm}{\sigma\sqrt n}
\xrightarrow{P}0.
\]

Значит при условиях ЦПТ закон больших чисел получается как более грубое
следствие: ЦПТ описывает форму флуктуаций, а ЗБЧ говорит только, что
среднее стабилизируется.

Главный пример - схема Бернулли. Если \(A\) происходит с вероятностью
\(p\), а

\[
X_k=\mathbf 1_A^{(k)},
\]

то

\[
\mathbb EX_k=p.
\]

Частота события равна

\[
\nu_n(A)=\frac{X_1+\cdots+X_n}{n}.
\]

По ЗБЧ:

\[
\nu_n(A)\xrightarrow{P}p=P(A).
\]

Это объясняет связь с аксиоматикой Колмогорова. В аксиоматике
вероятность - это мера на \(\sigma\)-алгебре событий, а не частота по
определению. Но закон больших чисел показывает, что в длинных сериях
независимых повторений относительные частоты стабилизируются около этой
меры. Поэтому аксиоматическая вероятность имеет частотную интерпретацию
и применима к экспериментам.

В конце важно добавить оговорку: не всякий ЗБЧ является частным случаем
ЦПТ, потому что ЗБЧ может быть верен при более слабых условиях. Но при
конечной дисперсии ЦПТ действительно влечет слабый закон больших чисел.

\subsection{15. Если осталось 2 минуты перед
экзаменом}\label{ux435ux441ux43bux438-ux43eux441ux442ux430ux43bux43eux441ux44c-2-ux43cux438ux43dux443ux442ux44b-ux43fux435ux440ux435ux434-ux44dux43aux437ux430ux43cux435ux43dux43eux43c}

Запомнить три формулы.

ЗБЧ:

\[
\overline X_n=\frac{S_n}{n}\xrightarrow{P}m.
\]

ЦПТ:

\[
\frac{S_n-nm}{\sigma\sqrt n}
\Rightarrow
N(0,1).
\]

Вывод:

\[
\overline X_n-m
=
\frac{\sigma}{\sqrt n}
\frac{S_n-nm}{\sigma\sqrt n}
\xrightarrow{P}0.
\]

Для события \(A\):

\[
\frac{N_n(A)}{n}\xrightarrow{P}P(A).
\]

Главная фраза: вероятность в аксиоматике Колмогорова задается мерой, а
закон больших чисел объясняет, почему эта мера наблюдается как
устойчивая частота в длинной серии независимых испытаний.

\newpage

\end{document}
